\Article{16}
In Art. 12---15 haben wir angenommen, dass die Fehler \(x,x',x''\) etc. zu Beobachtungen derselben Art gehören, sodass ihre einzelnen Wahrscheinlichkeiten durch dieselbe Funktion ausge\SourceLeaf{021}drückt werden. Es ist aber von selbst klar, dass sich die allgemeine Untersuchung von Art. 12---14 ebenso leicht auf den allgemeineren Fall erstreckt, in dem die Wahrscheinlichkeiten der Fehler \(x,x',x''\) etc. durch verschiedene Funktionen \(\varphi x,\varphi'x',\varphi''x''\) etc. ausgedrückt werden, das heißt, in dem jene Fehler zu Beobachtungen verschiedener Genauigkeit oder Unsicherheit gehören. Nehmen wir an, \(x\) sei der Fehler einer Beobachtung, deren zu befürchtender mittlerer Fehler \(=m\) ist; ferner seien \(x',x''\) etc. die Fehler anderer Beobachtungen, deren zu befürchtende mittlere Fehler jeweils \(m',m''\) etc. sind. Dann wird der Mittelwert der Summe \(xx+x'x'+x''x''+\text{etc.}\) gleich \(mm+m'm'+m''m''+\text{etc.}\) sein. Steht nun anderweitig fest, dass die Größen \(m,m',m''\) etc. in einem gegebenen Verhältnis stehen, etwa jeweils den Zahlen \(1,\mu',\mu''\) etc. proportional sind, so wird der Mittelwert des Ausdrucks
\[
\frac{xx+x'x'+x''x''+\text{etc.}}{1+\mu'\mu'+\mu''\mu''+\text{etc.}}
\]
\(=mm\) sein. Setzen wir aber einen bestimmten Wert desselben Ausdrucks, wie der Zufall die Fehler \(x,x',x''\) etc. liefert, gleich \(mm\), so ergibt sich der mittlere Fehler, dem diese Bestimmung weiterhin unterliegt, ähnlich wie im vorigen Artikel
\[
=\frac{\BareRoot{(n^4+n'{}^4+n''{}^4+\text{etc.}-m^4-m'{}^4-m''{}^4-\text{etc.})}}{1+\mu'\mu'+\mu''\mu''+\text{etc.}}
\]
wobei \(n',n''\) etc. in Bezug auf die Beobachtungen, zu denen die Fehler \(x',x''\) etc. gehören, dasselbe bedeuten sollen wie \(n\) in Bezug auf die erste Beobachtung. Wenn es also erlaubt ist, die Zahlen \(n,n',n''\) etc. den Größen \(m,m',m''\) etc. proportional anzunehmen, so wird jener zu befürchtende mittlere Fehler
\[
=\frac{\BareRoot{(n^4-m^4)}.\BareRoot{(1+\mu'{}^4+\mu''{}^4+\text{etc.})}}{1+\mu'\mu'+\mu''\mu''+\text{etc.}}
\]

Diese Art, einen Näherungswert von \(m\) zu bestimmen, ist aber nicht die zweckmäßigste. Um dies deutlicher zu zeigen, betrachten wir den allgemeineren Ausdruck
\[
y=\frac{xx+\alpha'x'x'+\alpha''x''x''+\text{etc.}}{1+\alpha'\mu'\mu'+\alpha''\mu''\mu''+\text{etc.}}
\]
dessen Mittelwert ebenfalls \(=mm\) sein wird, wie auch immer die Koeffizienten \(\alpha',\alpha''\) etc. gewählt werden. Der zu befürchtende mittlere Fehler, wenn wir einen bestimmten Wert von \(y\), wie der Zufall die Fehler \(x,x',x''\) etc. liefert, gleich \(mm\) annehmen, ergibt sich nach den oben angegebenen Grundsätzen
\[
=\frac{\BareRoot{(n^4-m^4+\alpha'\alpha'(n'{}^4-m'{}^4)+\alpha''\alpha''(n''{}^4-m''{}^4)+\text{etc.})}}{1+\alpha'\mu'\mu'+\alpha''\mu''\mu''+\text{etc.}}
\]
\SourceLeaf{022}Damit dieser mittlere Fehler möglichst klein wird, müssen wir setzen
\[
\begin{aligned}
\alpha'&=\frac{n^4-m^4}{n'{}^4-m'{}^4}.\mu'\mu'\\[4pt]
\alpha''&=\frac{n^4-m^4}{n''{}^4-m''{}^4}.\mu''\mu''\text{ etc.}
\end{aligned}
\]
Offensichtlich lassen sich diese Werte nicht berechnen, wenn nicht außerdem das Verhältnis der Größen \(n,n',n''\) etc. zu \(m,m',m''\) etc. anderweitig bekannt ist: Fehlt diese genaue Kenntnis, so scheint es wenigstens am sichersten\footnote{Wir verstehen nämlich die Kenntnis der Größen \(\mu',\mu''\) etc. nur in dem Fall als erreichbar, in dem nach der Natur der Sache die Fehler \(x,x',x''\) etc., die den Zahlen \(1,\mu',\mu''\) etc. proportional sind, als gleich wahrscheinlich anzusehen sind, oder vielmehr, in dem
\[
\varphi x=\mu'\varphi'(\mu'x)=\mu''\varphi''(\mu''x)\text{ etc.}
\]
}, jene diesen proportional anzunehmen (s. Art. 11), woraus sich die Werte
\[
\alpha'=\frac1{\mu'\mu'},\quad\alpha''=\frac1{\mu''\mu''}\text{ etc.}
\]
ergeben. Das heißt, die Koeffizienten \(\alpha',\alpha''\) etc. sind den relativen Gewichten der Beobachtungen gleichzusetzen, zu denen die Fehler \(x',x''\) etc. gehören, wobei das Gewicht der Beobachtung, zu der der Fehler \(x\) gehört, als Einheit genommen wird. Bezeichnet nun wie oben \(\sigma\) die Anzahl der vorgelegten Fehler, so erhalten wir den Mittelwert des Ausdrucks
\[
\frac{xx+\alpha'x'x'+\alpha''x''x''+\text{etc.}}{\sigma}
\]
\(=mm\) und den zu befürchtenden mittleren Fehler, wenn wir einen zufällig bestimmten Wert dieses Ausdrucks als den wahren Wert von \(mm\) annehmen,
\[
\frac{\BareRoot{(n^4+\alpha'\alpha'n'{}^4+\alpha''\alpha''n''{}^4+\text{etc.}-\sigma m^4)}}{\sigma}
\]
und folglich, wenn es tatsächlich erlaubt ist, die Größen \(n,n',n''\) etc. den Größen \(m,m',m''\) proportional anzunehmen,
\[
\mathord{\cdot}=\BareRoot{\frac{n^4-m^4}{\sigma}}
\]
Diese Formel ist mit derjenigen identisch, die wir oben für den Fall von Beobachtungen derselben Art gefunden haben.

\Article{17}
Ist der Wert einer Größe, die von einer anderen unbekannten Größe abhängt, durch eine Beobach\SourceLeaf{023}tung ohne absolute Genauigkeit bestimmt worden, so wird auch der daraus berechnete Wert der Unbekannten einem Fehler unterliegen; bei dieser Bestimmung bleibt aber nichts der Willkür überlassen. Sind dagegen \textit{mehrere} Größen, die von derselben Unbekannten abhängen, durch nicht absolut genaue Beobachtungen bekannt geworden, so können wir den Wert der Unbekannten entweder aus jeder einzelnen dieser Beobachtungen oder aus einer Kombination mehrerer Beobachtungen ermitteln, was auf unendlich viele verschiedene Arten möglich ist. Obwohl der so erhaltene Wert der Unbekannten immer einem Fehler unterworfen bleibt, ist doch bei der einen Kombination ein größerer, bei der anderen ein kleinerer Fehler zu befürchten. Ebenso verhält es sich, wenn mehrere Größen beobachtet worden sind, die zugleich von mehreren Unbekannten abhängen: Je nachdem, ob die Anzahl der Beobachtungen der Anzahl der Unbekannten gleich, kleiner oder größer ist, wird die Aufgabe bestimmt, unbestimmt oder mehr als bestimmt sein (wenigstens im Allgemeinen gesprochen); im dritten Fall lassen sich die Beobachtungen auf unendlich viele verschiedene Arten zur Bestimmung der Unbekannten kombinieren. Aus einer solchen Vielfalt von Kombinationen die zweckmäßigsten auszuwählen, das heißt diejenigen, die Werte der Unbekannten mit den kleinsten Fehlern liefern, ist gewiss die bei Weitem wichtigste Aufgabe bei der Anwendung der Mathematik auf die Naturphilosophie.

In der Theoria motus corporum coelestium haben wir gezeigt, wie die \textit{wahrscheinlichsten} Werte der Unbekannten zu ermitteln sind, wenn das Gesetz der Wahrscheinlichkeit der Beobachtungsfehler bekannt ist. Da dieses Gesetz seiner Natur nach in fast allen Fällen hypothetisch bleibt, haben wir jene Theorie auf das plausibelste Gesetz angewandt, bei dem die Wahrscheinlichkeit des Fehlers \(x\) der Exponentialgröße \(e^{-hhxx}\) proportional angenommen wird. Daraus ist die Methode hervorgegangen, die wir schon lange bei Rechnungen, besonders astronomischen, und jetzt auch die meisten Rechner unter dem Namen der Methode der kleinsten Quadrate verwenden.

Später hat der berühmte \textsc{Laplace} die Sache auf andere Weise angegriffen und gezeigt, dass dasselbe Prinzip noch immer allen anderen vorzuziehen ist, welches auch das Gesetz der Wahrscheinlichkeit der Fehler sei, sofern nur die Anzahl der Beobachtungen sehr groß ist. Bei einer mäßigen Anzahl von Beobachtungen blieb die Sache aber unberührt, sodass, wenn unser hypothetisches Gesetz verworfen wird, die Methode der kleinsten Quadrate nur deshalb vor anderen zu empfehlen ist, weil sie sich besonders für eine zweckmäßige Gestaltung der Rechnungen eignet.

Wir hoffen daher, dass es den Geometern willkommen sein wird, wenn wir bei dieser neuen Behandlung des Gegenstands zeigen, dass die Methode der kleinsten Quadrate die beste aller Kombinationen liefert. Und zwar nicht näherungsweise, sondern absolut, welches auch das Gesetz der Wahrscheinlichkeit der Fehler und welches die Anzahl der Beobachtungen sei, sofern wir nur den Begriff des \SourceLeaf{024}mittleren Fehlers nicht nach der Auffassung des berühmten \textsc{Laplace}, sondern so festlegen, wie wir es in Art. 5 und 6 getan haben.

Im Übrigen ist hier ausdrücklich vorauszuschicken, dass sich alle folgenden Untersuchungen nur auf unregelmäßige Fehler beziehen, die von einem konstanten Teil frei sind, da es zur vollkommenen Beobachtungskunst gehört, alle Ursachen konstanter Fehler mit größter Sorgfalt zu beseitigen. Welche Hilfsmittel der Rechner, der solche Beobachtungen zu behandeln unternimmt, bei denen ein begründeter Verdacht besteht, dass sie nicht frei von konstanten Fehlern sind, aus der Wahrscheinlichkeitsrechnung selbst schöpfen kann, behalten wir einer besonderen Untersuchung vor, die bei anderer Gelegenheit veröffentlicht werden soll.

\clearpage
